\section{The problem, stated precisely}
\label{sec:problem}

\begin{lede}
A warehouse floor, a never-ending stream of jobs, and a fleet that must keep
moving forever. This section fixes the model, the constraints and --- the part
that decides how every later number should be read --- the traffic regime the
measurements are taken in.
\end{lede}

\subsection{Environment and agents}

The warehouse is a 4-connected grid graph $G = (V, E)$. A vertex is a cell
$v = (r, c)$; $V$ holds the passable cells and $E$ joins orthogonally adjacent
ones. Shelving, walls and racking are simply absent from $V$. Certain vertices
carry a role: $V_{\mathrm{p}} \subset V$ are pickup stations,
$V_{\mathrm{d}} \subset V$ delivery stations, and $V_{\mathrm{k}} \subset V$
parking bays.

A fleet $A = \{1, \dots, n\}$ of identical robots occupies distinct vertices.
Time is discrete and synchronous: at each timestep every robot simultaneously
either moves to a vertex adjacent to its current one or stays where it is,
\begin{equation}
  \pos{i}{t+1} \in N(\pos{i}{t}) \cup \{\pos{i}{t}\},
  \label{eq:move}
\end{equation}
and a joint move is legal exactly when it creates neither of the two standard
conflicts:
\begin{align}
  \text{vertex conflict:} && \pos{i}{t+1} &= \pos{j}{t+1}, & i &\neq j,
  \label{eq:vertex-conflict}\\
  \text{swap conflict:} && \pos{i}{t+1} = \pos{j}{t} &\ \wedge\
    \pos{j}{t+1} = \pos{i}{t}, & i &\neq j.
  \label{eq:swap-conflict}
\end{align}
There is no rotation cost and no kinematic model; a robot's heading is tracked
only so that changing it can be discouraged (\Cref{sec:score-turn}).

\subsection{The task stream}

This is a \emph{lifelong} problem, not a one-shot one. Tasks are not known in
advance and never stop arriving. A task $\tau = (p, d, t_{\mathrm{rel}})$
releases at time $t_{\mathrm{rel}}$ and requires some robot to visit its pickup
$p \in V_{\mathrm{p}}$ and then its delivery $d \in V_{\mathrm{d}}$, in that
order. Arrivals follow one of four processes (\Cref{sec:arrivals}); the
experiments in this document use a Poisson process of mean rate $\lambda$
tasks per timestep.

A robot is assigned at most one task at a time and is free again on delivery.
Assignment is itself part of the problem --- which robot takes which job
matters --- and is handled by a greedy minimum-cost match
(\Cref{sec:assignment}) held \emph{identical} across every planner compared in
this document, so that comparisons isolate the movement layer.

The quantity a task contributes to the evaluation is its \emph{service time},
\begin{equation}
  \mathrm{svc}(\tau) = t_{\mathrm{complete}}(\tau) - t_{\mathrm{rel}}(\tau),
\end{equation}
defined only for tasks that actually complete --- a restriction
\Cref{sec:regime} returns to, because it is the single most misread number in
the whole evaluation.

\subsection{Objective}

Over a horizon of $T$ timesteps, with $K$ tasks completed, the primary
objective is \textbf{throughput}
\begin{equation}
  \Theta = \frac{K}{T} \quad \text{(tasks per timestep)},
  \label{eq:throughput}
\end{equation}
subject to the joint plan being conflict-free at every step. Makespan --- the
standard objective for one-shot MAPF --- is undefined here: there is no last
task. Sum-of-costs is likewise unavailable, since the set of costs is not
fixed in advance.

Secondary quantities are reported throughout, and each answers a different
question: mean and $p_{95}$ service time (latency, and tail latency),
Jain's fairness index over per-robot completion counts (is the fleet used
evenly), direction switches per 1000 steps (is the aisle signal stable),
head-on conflicts and counterflow moves (is the aisle layer doing what it
claims), and mean planner runtime per timestep (what the method costs).

\subsection{The traffic regime, and why it decides how to read every number}
\label{sec:regime}

The measurements in this document are taken in \textbf{saturation}, and this
must be stated before any of them is quoted.

The arrival rate $\lambda$ is set well above what the floor can serve. On
\variant{warehouse\_medium} tasks arrive at $\lambda = 1.5$ per timestep and
the best configuration completes about $0.50$ per timestep; on
\variant{warehouse\_bottleneck}, $\lambda = 0.8$ against $0.15$ served. Between
an eighth and a third of offered demand is served (\Cref{tab:scenarios}); the
backlog grows without bound for the whole run. Three consequences follow, and all three are
routinely misread.

\begin{enumerate}
\item \textbf{Throughput measures capacity, not effort.} In saturation there
is always another task waiting, so no planner is ever idle for want of work and
$\Theta$ is exactly the floor's service capacity under that planner. This is
why $\Theta$ is the primary objective and why differences in it are meaningful
even when they look small: $0.50$ against $0.31$ tasks per timestep is a $61\%$
difference in the capacity of the warehouse.

\item \textbf{Service time is censored, and the censoring is not random.}
$\mathrm{svc}(\tau)$ exists only for completed tasks. A planner that serves
only the nearby, easy tasks and abandons the rest reports an \emph{excellent}
mean service time. This is not hypothetical: on
\variant{warehouse\_bottleneck} Token Passing reports a mean service time of
$31$ timesteps --- against $154$ for the planner that delivers twenty times as
many tasks --- because the handful of tasks it finishes are the ones it
finishes in the first few timesteps, before the corridor jams for good. Mean
service time may be compared only between planners with comparable throughput.
\Cref{fig:censoring} shows this trap directly.

\item \textbf{Robot count is not a free parameter.} All planners are compared
at the same fleet size on the same map, because throughput is not monotone in
fleet size once the floor is congested --- adding robots past the congestion
knee lowers it.
\end{enumerate}

\begin{caveat}
Every table in \Cref{sec:results} therefore reports throughput first, quotes
service time only alongside it, and never ranks two planners by service time
alone. A reader who takes one thing from this section: \emph{a low mean service
time next to a low throughput is a symptom, not an achievement.}
\end{caveat}

\subsection{What makes this hard, in one paragraph}

A warehouse map is mostly corridors. In a single-file corridor, two robots
wanting opposite directions cannot both proceed, and no local rule that treats
them symmetrically resolves it; one must yield, and something must decide
which. Worse, in a lifelong setting the pair that meets head-on is replaced by
another pair a few timesteps later, so a resolution that works once must keep
working forever, under an adversarial arrival process the planner does not
control. Optimal joint planning would resolve this and is exponential
(\Cref{sec:related-cbs}); planning each robot independently against a
reservation table is cheap and deadlocks (\Cref{sec:related-tp}). The
interesting region is between the two, and that is where PIBT
(\Cref{sec:related-pibt}) and this work sit.
